\calculussheet{sheet:p18-erotetic-patterns}{Original IEL evocation and implication patterns}
 {Source \(L_1,L_2\) and native \(\supset,\equiv,?\pm|A,B|\).
 These are semantic E/Im warrants, not a complete primitive calculus.
 Every resulting question has at least two distinct closed answers.
 Families and notation: Definitions~\ref{def:p18-ep-notation}--\ref{def:p18-ep-families}.
 Source numbers are retained; proofs are supplied in the matched pattern annex.}

\hbrule{rule:p18-ep-e05}{E (6.5)}{p\lor\neg p}{?p}
 {Distinct propositional variables in E (6.5)--(6.15).
  All eleven proofs: Annex~\ref{annex:p18.ep.propositional-evocation}.}
\hbrule{rule:p18-ep-e06}{E (6.6)}{p\lor q}{?p}{\none}
\hbrule{rule:p18-ep-e07}{E (6.7)}{p\lor q}{?q}{\none}
\hbrule{rule:p18-ep-e08}{E (6.8)}{p\supset q}{?p}{\none}
\hbrule{rule:p18-ep-e09}{E (6.9)}{p\supset q}{?q}{\none}
\hbrule{rule:p18-ep-e10}{E (6.10)}{p\equiv q}{?\{p\land q,\neg p\land\neg q\}}{\none}
\hbrule{rule:p18-ep-e11}{E (6.11)}{p\lor q}{?\{p,q\}}{\none}
\hbrule{rule:p18-ep-e12}{E (6.12)}{p\lor q}
 {?\{p\land q,p\land\neg q,\neg p\land q\}}{\none}
\hbrule{rule:p18-ep-e13}{E (6.13)}{p\supset(q\lor r)\qquad p}{?\{q,r\}}{\none}
\hbrule{rule:p18-ep-e14}{E (6.14)}{(p\land q)\supset r\qquad\neg r}{?\{\neg p,\neg q\}}{\none}
\hbrule{rule:p18-ep-e15}{E (6.15)}{(p\lor q)\lor r}{?\{p,q\lor r\}}{\none}

\hbrule{rule:p18-ep-e16}{E (6.16)}{Pc_1,\ldots,Pc_n}{?\forall xPx}
 {Finite \(n\geq1\); all nonempty-domain \(L_2\)-models, not domain closure.
  Annex~\ref{annex:p18.ep.firstorder-polar}.}
\hbrule{rule:p18-ep-e17}{E (6.17)}{Pc_1\land Rc_1,\ldots,Pc_n\land Rc_n}
 {?\forall x(Px\supset Rx)}
 {Finite \(n\geq1\), distinct predicates. Annex~\ref{annex:p18.ep.firstorder-polar}.}
\hbrule{rule:p18-ep-e18}{E (6.18)}{Pc_1\land Rc_1,\ldots,Pc_n\land Rc_n,\ Pc_{n+1}}
 {?Rc_{n+1}}
 {Distinct constant symbols and predicates, no distinct-denotation axiom.
  Annex~\ref{annex:p18.ep.firstorder-polar}.}
\hbrule{rule:p18-ep-e19}{E (6.19)}{Pc_1,Pc_2}{?c_1=c_2}
 {Distinct symbols may denote the same individual. Annex~\ref{annex:p18.ep.firstorder-polar}.}
\hbrule{rule:p18-ep-e20}{E (6.20)}{\exists xPx}{?\{Pc,\exists x(Px\land x\ne c)\}}
 {Named-or-residual witness. Annex~\ref{annex:p18.ep.named-evocation}.}
\hbrule{rule:p18-ep-e21}{E (6.21)}{\exists x(Px\land(x=c_1\lor\cdots\lor x=c_n))}{?\{Pc_1,\ldots,Pc_n\}}
 {\(n>1\), distinct constant symbols. Annex~\ref{annex:p18.ep.named-evocation}.}
\hbrule{rule:p18-ep-e22}{E (6.22)}{\neg\forall x((x=c_1\lor\cdots\lor x=c_n)\supset Px)}
 {?\{\neg Pc_1,\ldots,\neg Pc_n\}}
 {\(n>1\), distinct constant symbols. Annex~\ref{annex:p18.ep.named-evocation}.}
\hbrule{rule:p18-ep-e23}{E (6.23)}{\exists x(Px\land(x=c_1\lor\cdots\lor x=c_n))}{?S(Px)}
 {\(n>1\); proof of non-entailment for every language constant.
  Annex~\ref{annex:p18.ep.named-evocation}.}
\hbrule{rule:p18-ep-e24}{E (6.24)}{\neg\forall x((x=c_1\lor\cdots\lor x=c_n)\supset Px)}{?S(\neg Px)}
 {\(n>1\); the whole wh-answer set, not just the finite negative list.
  Annex~\ref{annex:p18.ep.named-evocation}.}
\hbrule{rule:p18-ep-e25}{E (6.25)}
 {\forall x(Px\equiv(x=c_1\lor\cdots\lor x=c_n))\lor\cdots\lor
  \forall x(Px\equiv(x=c'_1\lor\cdots\lor x=c'_k))}{?U(Px)}
 {\(m\geq2\), \(n,k>1\), all intermediate alternative list lengths
  greater than one; distinct
  symbols across lists as in the source convention. Every candidate-list
  proof: Annex~\ref{annex:p18.ep.complete-list}.}

\hbrule{rule:p18-ep-im01}{Im (8.1)}{?\neg A}{?A}
 {Regular and analytic. Annexes~\ref{annex:p18.ep.pure-patterns},
  \ref{annex:p18.ep.analytic-profile}.}
\hbrule{rule:p18-ep-im02}{Im (8.2)}{?A}{?\neg A}
 {Regular and analytic. Same complete proofs.}
\hbrule{rule:p18-ep-im03}{Im (8.3)}{?\{A,B\}}{?\{B,A\}}
 {Distinct answers; regular and analytic. Same complete proofs.}
\hbrule{rule:p18-ep-im04}{Im (8.4)}{?\{A,B\lor C\}}{?\{A,B,C\}}
 {Arbitrary formulas; both questions well formed. Regular and analytic.
  Annex~\ref{annex:p18.ep.general-refinement}.}
\hbrule{rule:p18-ep-im05}{Im (8.5)}{?\{A,B,C\}}{?\{A,B\lor C\}}
 {Arbitrary formulas; both questions well formed. Analytic, not generally
  regular. Annex~\ref{annex:p18.ep.general-refinement}.}
\hbrule{rule:p18-ep-im07}{Im (8.7)}{?\pm|A,B|}{?(A\otimes B)}
 {\(\otimes\in\{\land,\lor,\supset,\equiv\}\).
  Analytic exactly for conjunction; not generally regular.
  Annexes~\ref{annex:p18.ep.pure-patterns},\ref{annex:p18.ep.analytic-profile}.}
\hbrule{rule:p18-ep-im10}{Im (8.10)}{?A}{?\pm|A,B|}
 {Regular; analytic iff both \(B,\neg B\in\mathcal H(?A)\).
  Same complete proofs; independent distinct atoms give non-analytic instances.}
\hbrule{rule:p18-ep-im11}{Im (8.11)}{?B}{?\pm|A,B|}
 {Regular; analytic iff both \(A,\neg A\in\mathcal H(?B)\).
  Same complete proofs.}
\hbrule{rule:p18-ep-im12}{Im (8.12)}{?(A\otimes B)\qquad\alpha A}{?B}
 {\(\alpha A\in\{A,\neg A\}\), fixed in each instance;
  all four connectives, regular. Annex~\ref{annex:p18.ep.signed-cofactors}.}
\hbrule{rule:p18-ep-im13}{Im (8.13)}{?(A\otimes B)\qquad\alpha B}{?A}
 {\(\alpha B\in\{B,\neg B\}\), fixed; all four connectives, regular.
  Annex~\ref{annex:p18.ep.signed-cofactors}.}
\hbrule{rule:p18-ep-im14}{Im (8.14)}{?\{\neg A,A\land B,A\land\neg B\}}{?A}
 {Empty basis even though printed in the nonempty-basis subsection.
  Not generally regular. Annex~\ref{annex:p18.ep.polar-refinements}.}
\hbrule{rule:p18-ep-im15}{Im (8.15)}{?\{A_1,\ldots,A_n\}\qquad A_1\lor\cdots\lor A_n}{?A_i}
 {\(n>1\), \(i\in J_n\); regular if \(n=2\).
  Annex~\ref{annex:p18.ep.polar-refinements}.}
\hbrule{rule:p18-ep-im16}{Im (8.16)}
 {?\{A_1,\ldots,A_n\}\qquad A_1\lor\cdots\lor A_n\qquad\neg(A_1\land\cdots\land A_n)}
 {?\{\neg A_1,\ldots,\neg A_n\}}
 {\(n>1\); regular if \(n=2\). Each complementary witness and both
  lost-premise cases: Annex~\ref{annex:p18.ep.denial}.}
\hbrule{rule:p18-ep-im18}{Im (8.18)}
 {?\{A_1,\ldots,A_n\}\qquad B\supset(A_1\lor\cdots\lor A_i)\qquad
  \neg B\supset(A_{i+1}\lor\cdots\lor A_n)}{?B}
 {\(1\leq i<n\); regular if \(n=2\). Annex~\ref{annex:p18.ep.partition-patterns}.}
\hbrule{rule:p18-ep-im19}{Im (8.19)}
 {?\{A_1,\ldots,A_n\}\qquad B\supset(A_1\lor\cdots\lor A_i)\qquad
  C\supset(A_{i+1}\lor\cdots\lor A_n)\qquad B\lor C}
 {?\{B,C\}}
 {\(1\leq i<n\); regular if \(n=2\). Same complete block proof.}
\hbrule{rule:p18-ep-im20}{Im (8.20)}
 {?\{A_1,\ldots,A_n\}\qquad(A_1\lor\cdots\lor A_i)\equiv B\qquad
  (A_{i+1}\lor\cdots\lor A_n)\equiv C}{?\{B,C\}}
 {\(1\leq i<n\); regular if \(n=2\). Same complete block proof.}
\hbrule{rule:p18-ep-im21}{Im (8.21)}
 {?\{A_1,\ldots,A_n\}\qquad(A_1\lor\cdots\lor A_i)\supset B\qquad
  (A_{i+1}\lor\cdots\lor A_n)\supset\neg B\qquad A_1\lor\cdots\lor A_n}{?B}
 {\(1\leq i<n\); regular if \(n=2\). The first displayed disjunction
  is the antecedent of its conditional, not an extra premise. Same proof.}
\hbrule{rule:p18-ep-im22}{Im (8.22)}
 {?A(x/c)\qquad\forall x(Ax\equiv(Bx\land Cx))}{?\pm|B(x/c),C(x/c)|}
 {Arbitrary \(Ax,Bx,Cx\) with only x free; closed substitution, no capture;
  regular. Annex~\ref{annex:p18.ep.arbitrary-formulas}.}
\hbrule{rule:p18-ep-im26}{Im (8.26)}
 {?S(Ax)\qquad\exists xAx\qquad\forall x(Ax\supset(x=c_1\lor\cdots\lor x=c_n))}{?A(x/c_i)}
 {\(n>1\), distinct constant symbols, \(i\in J_n\); arbitrary Ax with
  only x free. Positive singleton and negative complementary witnesses;
  regular if \(n=2\). Annex~\ref{annex:p18.ep.wh-polar}.}
